\documentclass[aps,pra,twocolumn,superscriptaddress,nofootinbib]{revtex4-2}

\usepackage{amsmath,amssymb,bm}
\usepackage{graphicx}
\usepackage{booktabs}
\usepackage{siunitx}
\usepackage{xcolor}
\usepackage{hyperref}

\hypersetup{colorlinks=true,citecolor=blue,urlcolor=blue,linkcolor=blue}
\sisetup{detect-all}

\newcommand{\Deff}{D_{\mathrm{eff}}}
\newcommand{\Dmicro}{D_{\mathrm{micro}}}
\newcommand{\gammac}{\gamma}
\newcommand{\Rregion}{\mathcal{R}}
\newcommand{\placeholderfigure}[1]{%
  \fbox{\parbox[c][1.45in][c]{0.94\columnwidth}{\centering\small #1}}}

\begin{document}

\title{Apparent diffusion from deterministic velocity gradients in wavelength-resolved unipolar photodetector transport}

\author{Anonymous}

\date{August 16, 2026}

\begin{abstract}
Wavelength-dependent photogeneration can encode internal transport in a detector response, but extracting a microscopic diffusion coefficient requires separating material transport from device-level heterogeneity sampled by finite generation profiles. We study this attribution problem for a single-mobile-carrier, or unipolar, planar Shockley--Ramo observable with zero microscopic diffusion and zero recombination. Finite generation-depth kernels that overlap a spatially nonuniform deterministic velocity region can yield a positive effective diffusion coefficient when the resulting terminal-current response is fit with a homogeneous drift--diffusion model. An exact full-contact planar continuum calculation gives $D_{\rm eff}=2.6182\times10^{-3}$, $2.5508\times10^{-3}$, and $2.3506\times10^{-3}\,\mathrm{m^2/s}$ at 100 MHz, 500 MHz, and 1 GHz, while upstream point-source controls remain at numerical-zero diffusion. Removing kernel support in the heterogeneous region collapses the effect even when every channel mean depth is restored. We derive the local low-frequency equivalence, an exact finite-support leakage relation, the upstream affine-plus-exponential bridge, and a profiled parameter-bias law separating tangent bias from normal model rejection. Independent velocity profiles give positive apparent diffusion for downstream acceleration, zero for uniform velocity, and negative values for deceleration. Root-space and direct full-channel multi-frequency tests quantify complementary falsification information: through 1 GHz the reference full-channel test lowers the required rejection SNR from 90.37 to 81.51 dB, whereas by 3 GHz the two protocols are comparable. Structured covariance and optical-kernel uncertainty stresses show that the numerical effective coefficient is conditional on the inverse metric and optical model. A separate pair-aware scope audit shows that the fitted single-carrier coefficient is not invariant when a second carrier contribution is admitted, so no generic two-carrier photodiode claim is made. The HgCdTe construction is retained only as a conditional optical/field/timing scale example, not as a self-consistent device simulation.
\end{abstract}

\maketitle

\section{Introduction}

The temporal response of a photodetector depends not only on microscopic carrier parameters but also on where photocarriers are generated, the electric field they encounter, and how their motion is weighted by the measurement electrodes. In high-speed photodiodes, finite absorption depth, depleted and undepleted transport regions, drift, diffusion, and device electrostatics have therefore long been treated together in forward response models \cite{hu2015pda}. Related work on photodiode optoelectronic chromatic dispersion (OED) has made the wavelength dependence of this internal transport response an observable in its own right: wavelength-dependent absorption changes the distribution of carrier transit and diffusion delays and consequently changes the radio-frequency (RF) phase and amplitude measured from a photodiode \cite{glasser2021oed,liokumovitch2021germanium,liokumovitch2022femtometer,dutta2024pin,mudgal2024pn}. Multi-frequency OED features have recently also been used as an inverse data space for spectral reconstruction \cite{kassa2026spectroscopy}.

Those established results motivate a different inverse question. Suppose wavelength is used not primarily to infer the optical spectrum, but to program a family of internal generation profiles from which a material transport parameter is inferred. When is a fitted diffusion coefficient attributable to microscopic diffusion rather than to deterministic device-level heterogeneity sampled differently by the finite optical kernels?

This distinction is not guaranteed by a good fit to a spatially simple transport model. Semiconductor time-of-flight analyses already show that inhomogeneous electric fields can bias transport extraction when a homogeneous-field model is imposed \cite{emelianova2006tof}, and electrostatic assumptions can bias diffusion-related quantities in other photocarrier inverse methods \cite{hattori1992sspg}. More generally, terminal current is an electrode-weighted observable rather than a direct readout of one microscopic carrier coordinate; consistent transient-current theory therefore requires the device electrostatics and observation operator to be distinguished from internal carrier motion \cite{hawks2015current}.

Here we isolate a specific material-attribution failure mode in a wavelength-resolved Shockley--Ramo measurement \cite{shockley1938,ramo1939}. The central construction is explicitly a single-mobile-carrier, or unipolar, terminal-current contribution; it is not a complete electron--hole-pair model of a generic photodiode transient. The forward systems used for the central counterexamples have
\begin{equation}
\Dmicro=0,
\qquad
\kappa=0,
\end{equation}
so all carrier trajectories are deterministic. Nevertheless, when finite generation kernels treated as known by the inverse overlap a downstream region of nonuniform deterministic velocity, fitting the resulting unipolar terminal-current response with a homogeneous drift--diffusion model returns
\begin{equation}
\Deff>0.
\end{equation}
The fitted parameter is mathematically valid as an effective parameter of the wrong homogeneous model; the error is interpreting it as microscopic material diffusion.

We establish four results. First, deterministic downstream acceleration generates the same low-frequency real spatial exponent that a homogeneous inverse associates with positive diffusion. Second, mean generation depth is not a sufficient optical coordinate: all nominal means may lie upstream of the nonuniform region while finite kernel support inside that region drives the bias. Removing that support collapses the apparent diffusion even if every channel mean is restored exactly. Third, the parameter bias and the same-frequency model residual are different projections of the channel discrepancy; a nuisance can therefore move the recovered root while producing little model-rejection signal. Fourth, additional RF frequencies reject the wrong model only when the covariance-weighted departure from the homogeneous dispersion manifold becomes statistically resolvable. This separates structural overdetermination from practical falsification.

The analysis is intentionally framed as a systematic-error and identifiability theory, not as a claim that wavelength-dependent RF photodiode response, field-dependent carrier transport, or effective-parameter bias are individually new phenomena. The contribution is the detector-specific combination of finite generation kernels supplied exactly to the inverse, deterministic velocity heterogeneity, Shockley--Ramo terminal-current inversion, a true-zero-diffusion control, causal support ablations, and quantitative attribution and rejection criteria.

Spectral-response transport extraction is itself established prior art, and optical-model error can bias such inferred diffusion quantities. In particular, Ashry and Fares reported strong sensitivity of photodiode diffusion-length extraction to absorption-coefficient errors \cite{ashry2003diffusion}. We therefore do not claim that optical-model sensitivity, generalized least squares, nuisance projection, or Schur-complement information is new. The detector-specific question developed below is the joint Shockley--Ramo spectral-depth attribution geometry under an explicit zero-microscopic-diffusion control.


\section{Measurement model and attribution problem}
\label{sec:model}

\subsection{Point-source terminal response}

Consider a planar absorber with coordinate
\begin{equation}
0\le z\le L,
\end{equation}
where the collecting electrode is at $z=L$. We use the Fourier convention $e^{-i\omega t}$. Let $H(z,\omega)$ denote the complex terminal-current transfer generated by a point source at depth $z$.

The observable $H$ below is one mobile-carrier contribution. In an ordinary electron--hole pair both carrier trajectories can contribute to terminal current; that more general decomposition is not assumed here. The present equations apply directly to a unipolar observable and to regimes in which the complementary carrier contribution is independently isolated or negligible. A concrete photodetector class in which strongly carrier-asymmetric transport is deliberately engineered is the \mbox{Uni-Traveling-Carrier photodiode} (UTC-PD), introduced to exploit electron transport in the collection region \cite{ishibashi1997utc}. We cite UTC-PDs only as an existence example for a physically meaningful effectively one-traveling-carrier regime; the planar surrogate below is not a UTC-PD device model.

For deterministic downstream velocity $v(z)>0$, zero microscopic diffusion, zero recombination, and the planar weighting potential $\phi_w=z/L$, the Shockley--Ramo path integral \cite{shockley1938,ramo1939} can be written
\begin{equation}
H(z,\omega)
=\frac{1}{L}\int_z^L
\exp\left[-i\omega\tau(x;z)\right]dx,
\label{eq:Hpath}
\end{equation}
with
\begin{equation}
\tau(x;z)=\int_z^x\frac{du}{v(u)}.
\end{equation}
Differentiation with respect to source coordinate gives the exact source-response equation
\begin{equation}
\boxed{
\frac{\partial H}{\partial z}
=-\frac{1}{L}+\frac{i\omega}{v(z)}H,
\qquad H(L,\omega)=0.
}
\label{eq:ramoode}
\end{equation}
The affine offset in $H$ means that the local spatial exponential is carried by the source derivative rather than by $H$ itself. Define
\begin{equation}
P(z,\omega)=\frac{\partial H}{\partial z}
\end{equation}
and the infinitesimal point-source spatial exponent
\begin{equation}
\boxed{
\gamma_{\rm loc}(z,\omega)
=-\frac{\partial}{\partial z}
\ln[-L P(z,\omega)].
}
\label{eq:gammaloc}
\end{equation}
For a finite source-coordinate step $h$, the corresponding local difference exponent is
\begin{equation}
\gamma_h(z,\omega)
=-\frac{1}{h}
\ln\!\left[\frac{P(z+h,\omega)}{P(z,\omega)}\right],
\end{equation}
which approaches Eq.~\eqref{eq:gammaloc} as $h\rightarrow0$.  This is the point-source limit of the first-difference exponential structure used by the spectral inverse.  The finite-kernel results below do not follow by evaluating $\gamma_{\rm loc}$ at a kernel mean; they require the separate kernel averaging in Eq.~\eqref{eq:channel}.

For constant velocity, the source-coordinate difference response is one exact spatial exponential. More strongly, suppose the source lies in a uniform upstream interval $z<z_d$ with velocity $v_0$, while arbitrary deterministic heterogeneity is allowed downstream of $z_d$. Equation~\eqref{eq:ramoode} then gives
\begin{equation}
\boxed{H(z,\omega)=\frac{v_0}{i\omega L}+C(\omega)e^{i\omega z/v_0},\qquad z<z_d.}
\label{eq:upstreambridge}
\end{equation}
All downstream heterogeneity enters only through the matching constant $C(\omega)$. Thus point sources confined to the uniform upstream interval remain exactly affine plus one source-coordinate exponential even though those carriers later traverse the heterogeneous region. A velocity gradient changes the source-coordinate law only when the source samples the nonuniform interval or a finite generation kernel has support there.

\subsection{Finite generation kernels treated as known by the inverse}

Let $g_m(z)$ be the normalized generation kernel of spectral channel $m$,
\begin{equation}
\int_0^L g_m(z)\,dz=1,
\end{equation}
and define its mean generation depth
\begin{equation}
\bar z_m=\int_0^L z g_m(z)\,dz.
\end{equation}
The measured complex channel is
\begin{equation}
\boxed{
J_m(\omega)=\int_0^L g_m(z)H(z,\omega)\,dz.
}
\label{eq:channel}
\end{equation}
We emphasize the distinction between the full kernel $g_m$ and the scalar mean $\bar z_m$. The former, not the latter, determines which device regions contribute to the channel.

The kernel-aware one-mode inverse is written in the numerically stable form
\begin{equation}
J_m=C+K F_m(r),
\label{eq:onemode}
\end{equation}
where
\begin{equation}
F_m(r)=\int_0^L g_m(z)
\frac{e^{r(z-z_{\rm ref})}-1}{r}\,dz
\end{equation}
with its continuous affine limit at $r=0$. The physical transport exponent used below is
\begin{equation}
\gamma=-r.
\end{equation}
This form retains the independently specified channel shapes rather than replacing them with delta functions or assuming rigid translations.  Throughout this work, a kernel treated as ``known'' means that the same theoretically specified $g_m(z)$ is supplied to the forward average and to the inverse.  No experimental kernel calibration is performed here, and optical-model uncertainty is deliberately excluded from the central counterexample.

For the conditional HgCdTe numerical example we use six normalized finite kernels generated by a monotonic graded-HgCdTe optical model using the Hansen band-gap relation and the Moazzami above-bandgap absorption parameterization \cite{hansen1982gap,moazzami2005absorption}.  The wavelength of each channel is chosen numerically so that the resulting normalized kernel has mean depth $2.0$, $2.5$, $3.0$, $3.5$, $4.0$, or $4.5\,\mu\mathrm{m}$ in the $7.6\,\mu\mathrm{m}$ absorber.  These theoretical kernels are used exactly in both the forward average and the kernel-aware inverse.  Their purpose is to provide a physically motivated finite-support family for the identifiability stress, not to represent measured optical kernels from a specific detector.  Uncertainty in the kernel model, optical interference, and experimental calibration error are not included in the central counterexample and are discussed as additional nuisances below.

\subsection{Homogeneous inverse model}

For homogeneous drift--diffusion--recombination under the present Fourier convention, the spatial exponent satisfies
\begin{equation}
D\gamma^2+w\gamma=\kappa-i\omega.
\label{eq:ddlaw}
\end{equation}
The central counterexample asks whether a heterogeneous deterministic forward response with $\Dmicro=0$ can nevertheless produce fitted roots that are well represented by Eq.~\eqref{eq:ddlaw} with $D=\Deff>0$.

\begin{figure}[t]
\centering
\includegraphics[width=\columnwidth]{numerics/paper02_figures/fig1_generation_kernels.pdf}
\caption{Measurement geometry and finite generation support. All six mean generation depths in the conditional HgCdTe stress lie upstream of the nonuniform region, but the theoretical kernels retain finite probability inside it.}
\label{fig:kernels}
\end{figure}

\section{Deterministic heterogeneity can appear diffusive}\label{sec:heterogeneity}
\label{sec:deterministic}

\subsection{Low-frequency equivalence}

Let the recovered spatial exponent be analytic near dc and write
\begin{equation}
\delta\gamma(\omega)
\equiv\gamma(\omega)-\gamma_0
=-ia_1\omega+a_2\omega^2+ia_3\omega^3+O(\omega^4).
\label{eq:gammaexpansion}
\end{equation}
For the homogeneous law, define
\begin{equation}
V_*=w+2D\gamma_0.
\end{equation}
After subtracting the dc equation, Eq.~\eqref{eq:ddlaw} becomes
\begin{equation}
D(\delta\gamma)^2+V_*\delta\gamma=-i\omega.
\end{equation}
The physical root has expansion
\begin{equation}
\delta\gamma_{\rm DD}
=-\frac{i\omega}{V_*}
+\frac{D\omega^2}{V_*^3}
+\frac{2iD^2\omega^3}{V_*^5}
+O(\omega^4).
\label{eq:ddexpansion}
\end{equation}
Consequently, for a locally physically admissible parameterization, a recovered response with $a_1>0$ and $a_2>0$ admits an effective homogeneous model matching it through quadratic order,
\begin{equation}
\boxed{
V_{*,\rm eff}=\frac{1}{a_1},
\qquad
\Deff=\frac{a_2}{a_1^3}.
}
\label{eq:effectiveD}
\end{equation}
The first locally non-adjustable coefficient is cubic: the homogeneous model predicts $a_3=2a_2^2/a_1$. Thus an additional low-frequency point is structurally overdetermining without necessarily being practically discriminating; a wrong mechanism can be tangent to the homogeneous model family through $O(\omega^2)$.

\subsection{Deterministic field-gradient contribution}

The connection to the local exponent is obtained by expanding the source derivative $P=\partial_zH$.  From Eq.~\eqref{eq:ramoode}, write
\begin{equation}
F(z)=\frac{L-z}{v(z)},
\qquad
I(z)=\int_z^L\frac{L-u}{v(u)}du.
\end{equation}
Then
\begin{equation}
-LP(z,\omega)
=1-i\omega F(z)-\omega^2\frac{I(z)}{v(z)}+O(\omega^3).
\end{equation}
Expanding $\ln[-LP]$ in Eq.~\eqref{eq:gammaloc} therefore gives the point-source coefficients quoted below.  This derivation applies to the local first-difference limit; finite optical kernels are treated separately in Sec.~\ref{sec:kernels}.

Expanding Eq.~\eqref{eq:Hpath} gives
\begin{equation}
H(z,\omega)=\frac{L-z}{L}
-\frac{i\omega}{L}I(z)+O(\omega^2),
\end{equation}
where
\begin{equation}
I(z)=\int_z^L\frac{L-u}{v(u)}du.
\end{equation}
For infinitesimally spaced point-source measurements, the quadratic coefficient of the apparent spatial exponent is
\begin{equation}
\boxed{
a_2(z)=
\frac{v'(z)}{v(z)^2}
\left[
\frac{(L-z)^2}{v(z)}-
\int_z^L\frac{L-u}{v(u)}du
\right].
}
\label{eq:a2field}
\end{equation}
If velocity increases monotonically downstream, $v'(z)>0$ and $v(u)\ge v(z)$ for $u\ge z$. Then
\begin{equation}
I(z)\le\frac{(L-z)^2}{2v(z)},
\end{equation}
so Eq.~\eqref{eq:a2field} gives $a_2(z)>0$. A deterministic accelerating carrier can therefore acquire a positive apparent homogeneous diffusion coefficient even though its microscopic trajectory contains no random walk.

In the weak-gradient limit,
\begin{equation}
\boxed{
\Deff(z)\simeq\frac{1}{2}(L-z)^2v'(z).
}
\label{eq:weakgradient}
\end{equation}
Equation~\eqref{eq:weakgradient} is an apparent-parameter relation defined by the chosen homogeneous inverse; it is not a replacement for microscopic diffusion physics.

\subsection{Independent profile test}

To separate this effect from any particular electrostatic solver, we solve Eq.~\eqref{eq:ramoode} for prescribed one-dimensional velocity profiles. The upstream region has common velocity $v_0$ and the downstream region follows either a linear or exponential profile whose endpoint ratio is $R=v(L)/v_0$. The optical kernels and inverse procedure are held fixed while $\Dmicro=0$ for every case.

The result is sign controlled. Uniform velocity gives $\Deff\simeq0$; every tested accelerating profile ($R>1$) gives $\Deff>0$; and every tested decelerating profile ($R<1$) gives $\Deff<0$. For example, $R=2$ yields $\Deff=3.30\times10^{-3}\,\mathrm{m^2/s}$ for the linear family and $2.69\times10^{-3}\,\mathrm{m^2/s}$ for the exponential family, whereas $R=0.5$ yields negative values of comparable scale. The deterministic full-channel relative residual remains of order $10^{-4}$ for the accelerating examples.  Because that normalization includes fitted offset and amplitude components, we do not treat it as a statistical acceptance criterion; the covariance-aware same-frequency test is given in Sec.~\ref{sec:bias}.

\begin{figure}[t]
\centering
\includegraphics[width=\columnwidth]{numerics/paper02_figures/fig2_apparent_diffusion_velocity_profiles.pdf}
\caption{Sign-controlled apparent diffusion from deterministic velocity heterogeneity. The forward model has $\Dmicro=0$ for every point. Uniform velocity gives $\Deff\simeq0$, acceleration gives $\Deff>0$, and deceleration gives $\Deff<0$.}
\label{fig:profiles}
\end{figure}

\section{Finite kernels couple mean-upstream generation to device heterogeneity}
\label{sec:kernels}

\subsection{Exact finite-support leakage}

Let a reference point-source response be
\begin{equation}
H_0(z)=A+B e^{rz}.
\end{equation}
For the downstream-heterogeneity construction of Sec.~\ref{sec:heterogeneity}, $H_0$ need not be the response of a globally uniform comparison device. At each frequency we may choose it as the analytic continuation of the \emph{actual} upstream solution in Eq.~\eqref{eq:upstreambridge}, namely $A=v_0/(i\omega L)$, $B=C(\omega)$, and $r=i\omega/v_0$. Then $H(z)=H_0(z)$ identically throughout the uniform upstream interval, and the discrepancy is supported only where the actual response departs from that continuation. With this choice, suppose the true response differs only in a nuisance region $\Rregion$,
\begin{equation}
H(z)=H_0(z)+\delta H(z),
\qquad
\delta H(z)=0\quad z\notin\Rregion.
\end{equation}
Then the kernel-averaged channel is exactly
\begin{equation}
J_m=A+B M_m(r)+E_m,
\end{equation}
with
\begin{equation}
\boxed{
E_m=\int_{\Rregion}g_m(z)\delta H(z)\,dz.
}
\label{eq:leakage}
\end{equation}
Define the restricted generation probability
\begin{equation}
p_{m,\Rregion}=\int_{\Rregion}g_m(z)\,dz.
\end{equation}
If $p_{m,\Rregion}=0$ for every channel, then $E_m=0$ exactly regardless of the complexity of the response inside $\Rregion$. For normalized nonnegative kernels,
\begin{equation}
\boxed{
|E_m|\le p_{m,\Rregion}H_{\Rregion},
\qquad
H_{\Rregion}=\sup_{z\in\Rregion}|\delta H(z)|.
}
\label{eq:leakbound}
\end{equation}
The mean generation depth does not appear in this zero-overlap condition.

\subsection{Point-source and mean-preserving controls}

The conditional HgCdTe stress uses $L=7.6\,\mu\mathrm{m}$ and a collector-side nonuniform region beginning at $z_d=4.6\,\mu\mathrm{m}$. The six nominal mean source coordinates span $2.0$--$4.5\,\mu\mathrm{m}$, so all means lie upstream.

Ideal point sources at those coordinates give
\begin{equation}
\Deff=1.7\times10^{-12}\,\mathrm{m^2/s},
\end{equation}
numerically zero on the relevant scale. Point sources placed inside the nonuniform region instead give $\Deff=4.87\times10^{-3}\,\mathrm{m^2/s}$.

The exact full-contact planar continuum calculation with the physical finite kernels gives at $100\,\mathrm{MHz}$
\begin{equation}
\Deff=2.6182\times10^{-3}\,\mathrm{m^2/s},
\label{eq:hgcdteD}
\end{equation}
and gives $2.5508\times10^{-3}$ and $2.3506\times10^{-3}\,\mathrm{m^2/s}$ at $500\,\mathrm{MHz}$ and $1\,\mathrm{GHz}$. These mesh-free values are the primary numerical results for the full-contact planar stress. The separate two-dimensional field/trajectory solver reproduces them within $0.320\%$, $0.087\%$, and $0.071\%$, respectively, and is retained as an independent numerical reproduction and extension path for nonplanar geometries.
Their probabilities inside the nonuniform region increase from approximately $0.60\%$ for the shallowest channel to $38.89\%$ for the deepest channel.

A direct causal ablation removes all generation support inside the nonuniform region. Because simple truncation also shifts the channel means, we apply a stronger control: after removing all support in the nonuniform region, the surviving upstream density of each channel is exponentially tilted until its original mean depth is restored. The maximum mean-depth error is $9.3\times10^{-15}\,\mu\mathrm{m}$, while the recovered diffusion collapses to
\begin{equation}
\Deff=-7.7\times10^{-8}\,\mathrm{m^2/s}\simeq0.
\end{equation}
The exponentially tilted zero-overlap profiles are mathematical causal controls, not a claim that those modified kernel shapes correspond to a realizable optical spectrum or material structure.  Their role is to hold the scalar mean fixed while removing all direct sampling of the nuisance region.
Thus the controls establish that mean generation depth alone is insufficient and that support inside the nuisance region is a causal contributor to the material-attribution bias.  The exact leakage relation in Eq.~\eqref{eq:leakage} further shows that the detailed kernel shape within the nuisance region can matter in addition to the total overlap probability.

\begin{figure}[t]
\centering
\includegraphics[width=\columnwidth]{numerics/paper02_figures/fig3_tail_ablation.pdf}\\[2pt]%
\includegraphics[width=\columnwidth]{numerics/paper02_figures/fig3b_remote_overlap.pdf}
\caption{Generation support inside the nonuniform region is the causal optical variable. Scaling only the generation weight inside the nonuniform region produces a monotonic change in $\Deff$. Removing that support collapses the result even when every original mean generation depth is restored.}
\label{fig:ablation}
\end{figure}

\section{Parameter bias and model rejection}
\label{sec:bias}

\subsection{Profiled root-bias law}

At a fixed RF frequency, collect the specified finite-kernel channels into
\begin{equation}
\bm y=\bm f(C,K,r)+\bm E+\bm n,
\end{equation}
with
\begin{equation}
\bm f(C,K,r)=C\bm 1+K\bm F(r).
\end{equation}
Let $W$ be the positive-definite measurement weight. Define
\begin{equation}
X=\begin{bmatrix}\bm 1 & \bm F\end{bmatrix}
\end{equation}
and its weighted projector
\begin{equation}
P_X=X(X^\dagger W X)^{-1}X^\dagger W.
\end{equation}
The raw root-sensitivity direction is $\bm h=K\bm F_r$. After profiling the linear offset and amplitude,
\begin{equation}
\bm h_\perp=(I-P_X)\bm h.
\end{equation}
The first-order fitted-root bias is
\begin{equation}
\boxed{
\delta r=
\frac{\bm h_\perp^\dagger W\bm E}
{\bm h_\perp^\dagger W\bm h_\perp}.
}
\label{eq:rootbias}
\end{equation}
Equation~\eqref{eq:rootbias} makes explicit why parameter bias and model residual need not track one another. The component of $\bm E$ tangent to the kernel-aware one-mode manifold moves fitted parameters; only its normal component appears as first-order goodness-of-fit error.

Combining Eq.~\eqref{eq:rootbias} with the finite-support bound gives
\begin{equation}
\boxed{
|\delta r|
\le H_{\Rregion}
\frac{\left(\sum_m w_m p_{m,\Rregion}^2\right)^{1/2}}
{\|\bm h_\perp\|_W}
}
\label{eq:rootbound}
\end{equation}
for diagonal weights, to first order in the nuisance response.

Near a pure-drift baseline, $\gamma=a+ib$ has $|a|\ll|b|$ and $b\simeq-\omega/w$. Solving Eq.~\eqref{eq:ddlaw} for $D$ at $\kappa=0$ gives
\begin{equation}
D=-\frac{\omega a}{b(a^2+b^2)},
\end{equation}
so
\begin{equation}
\boxed{
D_{\rm app}\simeq\frac{w^3}{\omega^2}\Re\gamma.
}
\label{eq:driftsusceptibility}
\end{equation}
A small nuisance-induced real spatial exponent can therefore be interpreted as a substantial diffusion coefficient at low RF.

\subsection{Numerical validation of the bias law}

We validate Eq.~\eqref{eq:rootbias} at the exact uniform, zero-diffusion baseline using small positive and negative linear and exponential velocity perturbations while keeping the six theoretical optical kernels fixed. For $|\epsilon|\le0.002$, the projected first-order expression predicts the independently refitted complex root shift with maximum relative error $2.65\times10^{-6}$ and the propagated diffusion bias with maximum relative error $3.53\times10^{-4}$. For $|\epsilon|\le0.01$, the diffusion linearization error remains $1.77\times10^{-3}$.

Only about $3.44\times10^{-3}$ of the weak-gradient nuisance-vector amplitude norm lies outside the local kernel-aware one-mode tangent in this test, corresponding to a normal residual energy fraction of order $1.2\times10^{-5}$. The nuisance is therefore capable of moving the fitted material parameter far more efficiently than it produces a same-frequency rejection signal.

\begin{figure}[t]
\centering
\includegraphics[width=\columnwidth]{numerics/paper02_figures/fig4_bias_law_validation.pdf}
\caption{Validation of the local parameter-bias law. The first-order tangent-space prediction closely follows the nonlinear kernel-aware inverse for both independent velocity families.}
\label{fig:biasvalidation}
\end{figure}

\subsection{Same-frequency model rejection versus positive-$D$ detection}

A small deterministic fit residual relative to the full channel vector is not by itself a statistical model-acceptance statement. We therefore place the same-frequency one-mode goodness-of-fit test and the positive-$D$ test on the same explicit theoretical noise scale. At one RF frequency the six complex channels provide 12 real observations, while the complex $(C,K,r)$ one-mode model has six real fitted parameters, leaving six local residual degrees of freedom.

We retain the independent equal real/imaginary quadrature-noise model used below and define $S=\sqrt{\langle |J_m|^2\rangle}/\sigma_{\rm q}$. For the one-mode test, the deterministic post-fit channel residual supplies the noncentral alternative. For positive-$D$ detectability, the full $12\times6$ channel Jacobian is propagated through the fitted complex root and then through $D(\gamma)$; an idealized one-sided Gaussian $D>0$ test is evaluated at the same $\alpha=0.0027$ and 90\% power.

\begin{table}[t]
\caption{Same-frequency statistical ordering under the reference independent-quadrature noise model. $S_D$ is the RMS-channel SNR required for 90\% power to establish positive apparent diffusion; $S_{1\mathrm{m}}$ is the SNR required for 90\% power to reject the six-channel one-mode manifold, both at $\alpha=0.0027$.}
\label{tab:samefreq}
\centering
\begin{tabular}{rrrr}
\toprule
RF & $D_{\rm eff}$ (m$^2$/s) & $S_D$ (dB) & $S_{1\mathrm{m}}$ (dB)\\
\midrule
100 MHz & 2.610e-03 & 111.86 & 102.29\\
500 MHz & 2.549e-03 & 70.04 & 88.19\\
1000 MHz & 2.349e-03 & 52.41 & 81.80\\
\bottomrule
\end{tabular}
\end{table}

Under this reference covariance, the ordering is frequency dependent.  At the tested point(s) 100 MHz: $S_{1\mathrm{m}}=102.29$ dB before $S_D=111.86$ dB, so the spectral-model check self-announces before positive apparent diffusion reaches the stated 90\% detection power.  At the remaining tested point(s), positive apparent diffusion reaches 90\% power first, leaving finite RMS-channel-SNR windows (500 MHz: $70.04$--$88.19$ dB; 1 GHz: $52.41$--$81.80$ dB) in which $D>0$ has reached the stated detection power while rejection of the same-frequency one-mode manifold has not.  The example therefore supports conditional same-frequency hidden-risk windows at those RF points, not a universal stealth claim.

This comparison also clarifies the role of the deterministic full-vector residual quoted elsewhere in the paper: it is a descriptive approximation error, whereas Table~\ref{tab:samefreq} is the corresponding covariance-aware model-selection statement for the reference noise model.


\subsection{Root-space and full-channel multi-frequency rejection}

Additional RF frequencies test whether the recovered roots belong to the homogeneous dispersion family. Let the real-stacked multi-frequency root vector be $\bm x$, its covariance be $C_\gamma$, and the homogeneous model prediction be $\bm m(\bm p)$. The generalized least-squares rejection statistic is
\begin{equation}
\boxed{
T=\min_{\bm p}
[\bm x-\bm m(\bm p)]^T
C_\gamma^{-1}
[\bm x-\bm m(\bm p)].
}
\label{eq:teststat}
\end{equation}
Let $\bm d$ be the deterministic discrepancy from the closest homogeneous-model point and $J_p$ the model Jacobian. The covariance-weighted precision projected normal to the fitted-model tangent is
\begin{equation}
Q_\perp=C_\gamma^{-1}
-C_\gamma^{-1}J_p
(J_p^TC_\gamma^{-1}J_p)^{-1}
J_p^TC_\gamma^{-1},
\end{equation}
with alternative noncentrality
\begin{equation}
\boxed{
\Lambda=\bm d^TQ_\perp\bm d.
}
\label{eq:lambda}
\end{equation}
Practical discrimination is governed by $\Lambda$, not by the raw percent residual of one later frequency.

To illustrate the scale, we adopt an explicit theoretical noise model: six complex channels per RF frequency, independent equal Gaussian real and imaginary quadrature noise, equal RMS-channel SNR at each included RF, and no cross-frequency correlation. We define $S=\sqrt{\langle|J_m|^2\rangle}/\sigma_q$ and quote $S_{\rm dB}=20\log_{10}S$. In the root-space protocol the channel covariance is propagated through the kernel-aware one-mode root fit. At every cumulative RF bandwidth, the wrong homogeneous $D,w$ model is then re-fit jointly before its noncentral distance is evaluated. For $n$ frequencies, root space contains $2n$ real root coordinates and fits two real global parameters $(D,w)$, so its residual degrees of freedom are $\nu_{\rm root}=2n-2$. The full-channel statistic contains $12n$ real quadratures and profiles two complex nuisance coefficients $(C_f,K_f)$ at each frequency in addition to the same two real global parameters, giving $\nu_{\rm full}=12n-(4n+2)=8n-2$. Under the deterministic alternative the corresponding quadratic form is treated as noncentral $\chi^2_{\nu}(\Lambda)$ with the same residual dimension; power is evaluated from $P[\chi^2_{\nu}(\Lambda)>\chi^2_{\nu,1-\alpha}]$.

For false-rejection probability $\alpha=0.0027$ and $90\%$ power, the exact-continuum root-space test requires $90.37\,\mathrm{dB}$ through $1\,\mathrm{GHz}$, $79.90\,\mathrm{dB}$ through $1.5\,\mathrm{GHz}$, $73.20\,\mathrm{dB}$ through $2\,\mathrm{GHz}$, and $64.21\,\mathrm{dB}$ through $3\,\mathrm{GHz}$. A direct full-channel test, which retains same-frequency normal residuals while profiling complex $C_f,K_f$ at each frequency, requires $81.51\,\mathrm{dB}$ through $1\,\mathrm{GHz}$, $72.28\,\mathrm{dB}$ through $2\,\mathrm{GHz}$, and $65.00\,\mathrm{dB}$ through $3\,\mathrm{GHz}$. Thus full channel is substantially stronger at intermediate bandwidth, whereas by 3 GHz the lower-dimensional root-space statistic is slightly stronger for this particular alternative after its smaller residual degrees-of-freedom penalty. This is a test-specific power comparison: the 3-GHz reversal does not imply that root compression retains more raw information. Neither test is claimed globally optimal outside the declared model and covariance.

The wrong model does not fail abruptly. Its best-fit diffusion decreases from $2.61\times10^{-3}\,\mathrm{m^2/s}$ over the lowest band to approximately $1.03\times10^{-3}\,\mathrm{m^2/s}$ when the fit extends through $3\,\mathrm{GHz}$. The homogeneous model continuously sacrifices parameter consistency to remain near the heterogeneous response. The profiled manifold distance in Eq.~\eqref{eq:teststat} is therefore the relevant rejection quantity.

\begin{figure}[t]
\centering
\includegraphics[width=\columnwidth]{numerics/paper02_figures/fig5_required_snr_vs_bandwidth_rev9.pdf}\\[2pt]%
\includegraphics[width=\columnwidth]{numerics/paper02_figures/fig5b_refitted_diffusion_vs_bandwidth.pdf}
\caption{Multi-frequency discrimination of the deterministic nuisance from homogeneous diffusion. Top: required RMS-channel SNR for the root-space and direct full-channel tests under the stated $\alpha$ and power criterion. Full-channel retention of same-frequency normal residuals is advantageous at intermediate bandwidth; the small reversal at 3 GHz is test-specific and reflects the smaller root-space residual dimension rather than greater raw information after compression. Bottom: the homogeneous $D$ re-fitted in root space decreases as the usable bandwidth grows.}
\label{fig:snr}
\end{figure}


\section{Measurement covariance and optical-model uncertainty}
\label{sec:modeluncertainty}

The central counterexample deliberately supplies the same theoretical kernels to the forward average and inverse so that deterministic velocity heterogeneity can be isolated. Here we relax the measurement metric and optical-model assumptions separately. The projection identities below are standard generalized-least-squares and nuisance-parameter geometry and are used as attribution tools rather than claimed as new mathematics. Spectral photodiode inverses are already known to depend sensitively on optical-model inputs: Ashry and Fares extracted diffusion length from wavelength-dependent photodiode response and reported strong sensitivity to absorption-coefficient errors \cite{ashry2003diffusion}. The narrower issue here is whether an optical-model error can align with the same-frequency transport-root tangent of a Shockley--Ramo RF inverse.

For a real-stacked data vector with small deterministic discrepancy $\bm e$, fitted-model Jacobian $G$, and precision $W=\Sigma^{-1}$, the local generalized-least-squares displacement is
\begin{equation}
\boxed{\delta\bm\theta=(G^T W G)^{-1}G^T W\bm e},
\label{eq:generalbias}
\end{equation}
with weighted tangent projector $P_W=G(G^T W G)^{-1}G^T W$ and post-fit discrepancy $(I-P_W)\bm e$. Thus covariance changes both the parameter-bias direction and the normal distance used for rejection. Under model misspecification the pseudo-true parameter $\bm\theta_*(W)=\arg\min_{\bm\theta}\|\bm y-\bm f(\bm\theta)\|_W^2$ can itself depend on the metric.

We re-fit the exact-known-kernel heterogeneous response under twelve normalized covariance shapes: IID; channel equicorrelation and AR(1) through $\rho=0.8$; low-rank common, spectral-slope, and spectral-curvature directions; and real--imaginary quadrature correlation. The frequency-dependent ordering is unchanged throughout this tested family. At $100\,\mathrm{MHz}$ one-mode rejection precedes positive-$D$ detection in all 12 cases, whereas at $500\,\mathrm{MHz}$ and $1\,\mathrm{GHz}$ positive $D$ is detectable first in all 12. The required positive-$D$ versus rejection SNR spans are $104.9$--$121.8$ versus $95.3$--$102.3\,\mathrm{dB}$ at 100 MHz, $63.1$--$79.7$ versus $81.2$--$88.2\,\mathrm{dB}$ at 500 MHz, and $45.4$--$61.1$ versus $74.8$--$82.3\,\mathrm{dB}$ at 1 GHz. The fitted $D_{\rm eff}$ at 100 MHz spans $1.66\times10^{-3}$--$2.61\times10^{-3}\,\mathrm{m^2/s}$ across the same weighting metrics. A separate cross-frequency stress changes the 3-GHz rejection requirement by at most about $1.1\,\mathrm{dB}$ over the tested correlation family; the bandwidth advantage therefore survives these particular correlations.

Fixed optical-kernel error is a distinct nuisance. If the nominal kernels depend on optical coordinates $\bm\alpha$, then a small fixed error gives $\delta\bm J=B\delta\bm\alpha$, with
\begin{equation}
B_{mj}=\int_0^L\frac{\partial g_m(z;\bm\alpha)}{\partial\alpha_j}H(z)\,dz.
\end{equation}
Equation~\eqref{eq:generalbias} then separates its tangent bias from its normal rejection signal. A zero-mean random kernel uncertainty may contribute $BC_\alpha B^T$ to an effective covariance at linear order, but a fixed or biased calibration error retains its deterministic tangent component.

An exact affine-depth control distinguishes a benign coordinate rescaling from non-affine kernel error. In the uniform deterministic $D_{\rm micro}=0$ problem, $z_{\rm true}=a+bz$ gives $\gamma_{\rm eff}=b\gamma$, $D_{\rm eff}=D/b^2$, and $w_{\rm eff}=w/b$. An exact continuum calculation with a 1\% depth-scale compression moves kernel means by as much as $18\,\mathrm{nm}$ yet leaves $|D_{\rm eff}|<4.7\times10^{-14}\,\mathrm{m^2/s}$ through 1 GHz.

Non-affine channel-dependent kernel error can behave differently. As a controlled sensitivity coordinate, take $\delta\lambda_m=A(-1,-0.6,-0.2,0.2,0.6,1)_m$. In the exact uniform-velocity $D_{\rm micro}=0$ null, while the inverse still uses the nominal kernels, $A=0.02997\,\mathrm{nm}$ produces $D_{\rm eff}(100\,\mathrm{MHz})=2.6182\times10^{-3}\,\mathrm{m^2/s}$, matching the central heterogeneous exact-continuum value. The maximum kernel-mean shift is $0.206\,\mathrm{nm}$, the one-mode relative residual is $3.1\times10^{-8}$, and positive $D$ reaches the stated detection power at $115.2\,\mathrm{dB}$ while same-frequency rejection requires $156.8\,\mathrm{dB}$. A signed curvature-shaped nuisance reproduces the same target at $A=-0.00294\,\mathrm{nm}$ with a maximum mean-depth shift of $0.020\,\mathrm{nm}$ and thresholds $115.2$ and $138.1\,\mathrm{dB}$.

These sub-nanometer amplitudes are not instrument calibration specifications or empirical error bars. They are coordinates in particular theoretical perturbations of the wavelength-to-generation-kernel map and expose conditioning and tangent alignment. The exact-known-kernel assumption is therefore load-bearing for interpreting the numerical magnitude of $D_{\rm eff}$ as a material parameter, although the deterministic velocity-heterogeneity counterexample remains independently established under that assumption.

\section{HgCdTe scale example and implications}
\label{sec:hgcdte}

The preceding results do not depend on HgCdTe specifically.  This section checks only whether the effective-field and timing scales used by the conditional stress have independent precedent in HgCdTe; it does not validate the modeled velocity profile or the predicted $D_{\rm eff}$ for any published detector.  We use the existing theoretical HgCdTe kernel construction as a scale example because independently published graded HgCdTe studies provide relevant field and timing benchmarks.

The planar stress uses $L=7.6\,\mu\mathrm{m}$, fixed terminal bias $V_{\rm bias}=0.30\,\mathrm{V}$, and a collector-side Poisson-curvature region of width $W_d=3.0\,\mu\mathrm{m}$ with parameter $\Delta=0.05\,\mathrm{V}$, defined in the full-contact planar limit by $V''=2\Delta/W_d^2$.  Thus $\Delta/W_d=166.7\,\mathrm{V/cm}$ is a characteristic curvature-field scale, not an independently added terminal field.  With the endpoint potentials held fixed, the exact planar field magnitude spans approximately $328.9$--$662.3\,\mathrm{V/cm}$ across the curved region.  Its regional mean is $495.6\,\mathrm{V/cm}$, compared with $394.7\,\mathrm{V/cm}$ for the corresponding uniform-bias profile, so the mean field increment is $100.9\,\mathrm{V/cm}$. A published compositionally graded HgCdTe study reports calculated composition-gradient built-in fields of approximately $100$--$200\,\mathrm{V/cm}$ in the linear graded region, with stronger local surface fields when a nonlinear composition region is retained \cite{xu2023graded}. Separately, a graded-bandgap room-temperature HgCdTe photodetector has demonstrated a response near $1.33\,\mathrm{ns}$, corresponding to a frequency response around $750\,\mathrm{MHz}$, with the faster response attributed to the composition-gradient built-in field \cite{sang2022graded}.

These comparisons do not calibrate the conditional stress to either published detector. The optical kernels, exact field profiles, junction electrostatics, readout covariance, and excitation conditions differ.  The comparison is only one of scale: the exact mean field increment of $100.9\,\mathrm{V/cm}$ in the planar stress lies within the cited $100$--$200\,\mathrm{V/cm}$ composition-gradient range, while the modeled profile itself is not taken from the published device.  The literature therefore supports physical plausibility of the field scale, not validation of the predicted $D_{\rm eff}$ for a real detector.

The scale comparison also highlights the practical attribution problem. The physical heterogeneity can be relevant in a detector family whose usable response bandwidth remains in a regime where the wrong homogeneous model is difficult to reject under the reference covariance. Independent electrostatic constraints can therefore be as important as increasing low-frequency measurement precision.

More generally, the design implication is not that every spectral-depth experiment requires zero overlap with all nonuniform regions. Rather, the full supplied generation kernels should be evaluated against known device-level nuisance regions. If a bound on the point-response discrepancy $H_{\Rregion}$ is available, Eq.~\eqref{eq:rootbound} converts kernel exposure and inverse conditioning into a systematic root-bias bound. A microscopic diffusion attribution should exceed that bound and should remain consistent with a multi-frequency dispersion test whose covariance-weighted power is explicitly quantified.

\section{Discussion}\label{sec:discussion}

The present result is an inverse-model statement. Deterministic velocity heterogeneity is real transport physics; the word ``apparent'' refers to assigning its terminal-current signature to a homogeneous microscopic diffusion coefficient. That distinction matters because many practical detector models intentionally use effective parameters. If the scientific objective is only to compress a device response into a compact empirical model, $\Deff$ may be useful. If the objective is to infer a material diffusion coefficient, the same fitted value can be misleading.

Three aspects make the attribution problem particularly severe. First, supplying the exact theoretical optical kernels removes optical-model uncertainty from the controlled inverse but cannot make the underlying point-source transport homogeneous. Exact knowledge of $g_m(z)$ therefore does not by itself eliminate Eq.~\eqref{eq:leakage}. Second, same-frequency model residual and parameter bias probe different directions in channel space. A nuisance nearly tangent to the kernel-aware model can strongly move the root while producing a very small fit residual. Third, the low-frequency homogeneous diffusion family itself has enough freedom to reproduce the first two frequency coefficients of a different analytic response. Practical rejection therefore begins with higher-order frequency structure and can require substantially more bandwidth than an identification calculation alone would suggest.

The Supplemental Material specifies the theoretical optical-kernel construction, exact planar deterministic transport stress, independent two-dimensional numerical reproduction, pair-aware scope audit, and executable reproduction files.

For the full-width contact, the central calculation is evaluated directly in the exact planar continuum. It gives $D_{\rm eff}=2.6182\times10^{-3}$, $2.5508\times10^{-3}$, and $2.3506\times10^{-3}\,\mathrm{m^2/s}$ at $100\,\mathrm{MHz}$, $500\,\mathrm{MHz}$, and $1\,\mathrm{GHz}$. The upstream point-source sequence remains at numerical-zero scale, while the inside-region control remains positive. A separate two-dimensional field/trajectory calculation reproduces these values within $0.320\%$, $0.087\%$, and $0.071\%$ and is retained as an independent solver check rather than the definition of the planar result.

A pair-aware scope audit confirms why the unipolar restriction is necessary. The exact pair model satisfies $H_{\rm down}(z,0)+H_{\rm up}(z,0)=1$ to $1.1\times10^{-16}$ and its uniform two-mode null gives zero downstream diffusion to numerical precision. With a heterogeneous downstream carrier, however, a free two-root decomposition becomes underconstrained; even when the simple countercarrier propagation root is fixed to its known value, only one of 21 core speed/frequency cases retains positive downstream $D$. The fitted single-carrier coefficient is therefore not invariant under a generic two-carrier decomposition, and no such generalization is claimed.

The framework has several limitations. Most importantly, the theorem concerns a single-mobile-carrier/unipolar terminal-current contribution; a generic two-carrier transient introduces a separate decomposition problem. The main counterexample deliberately omits microscopic diffusion so its deterministic source is unambiguous, while the optical-model extension tests controlled nuisance directions rather than a validated experimental error distribution. Its sub-nanometer wavelength amplitudes therefore quantify conditioning of this surrogate and must not be read as instrument specifications. A real detector can contain true diffusion, deterministic heterogeneity, optical-model error, parasitics, and correlated calibration errors simultaneously. The local projection formulas assume a regular identified one-mode solution, and the covariance study spans broad structured families rather than arbitrary covariance. Full electrostatic-model uncertainty and optimized measurement-time allocation remain outside the present scope.

These limitations suggest direct extensions. A nuisance-aware joint model could include parameterized electrostatic profiles and compare their Fisher directions with those of microscopic diffusion. Wavelengths could be selected not only for depth separation but also to minimize restricted overlap with poorly characterized device regions or to maximize the normal distance between competing transport models. Multi-frequency experimental design could likewise allocate information where the heterogeneous and homogeneous dispersion manifolds diverge most strongly rather than uniformly increasing precision at all frequencies.

\section{Conclusion}

For a single-mobile-carrier or unipolar wavelength-resolved terminal-current observable, a positive fitted diffusion coefficient need not be a unique signature of microscopic diffusion. In deterministic zero-diffusion models, finite generation kernels that overlap a region of spatially varying carrier velocity can produce a terminal-current response that a homogeneous drift--diffusion inverse interprets as $\Deff>0$.  At the point-source level, monotonic downstream acceleration gives a positive quadratic apparent-diffusion coefficient; in the tested linear and exponential finite-kernel families the fitted sign follows the velocity-gradient direction.  The positive accelerating cases survive exact finite-kernel fitting and can remain close to the homogeneous low-RF dispersion law over a finite bandwidth.

The decisive optical variable is finite support in the heterogeneous region, not the nominal mean generation depth. Removing that support collapses the apparent diffusion even when the channel means are preserved. Locally, the nuisance-induced root bias is determined by its projection onto the kernel-aware model tangent, while the model-rejection residual is the complementary normal projection. This separation explains how an apparently excellent one-mode fit can coexist with a materially wrong diffusion attribution.

The optical-model stress adds a stricter attribution condition. A fitted diffusion magnitude can be interpreted as material-like only to the extent that kernel-nuisance directions overlapping the transport-root tangent are independently constrained. This does not replace the deterministic velocity-gradient counterexample; it bounds what its fitted effective coefficient can mean when the optical model is uncertain.


Additional RF frequencies provide a falsification route, but structural overdetermination is not equivalent to statistical power. Root-space and full-channel tests show that bandwidth and retained channel-space residual directions provide complementary nuisance discrimination; neither compression is uniformly superior over the entire tested bandwidth. Together, the support bound, parameter-bias law, and multi-frequency rejection criterion provide a framework for separating microscopic material diffusion from deterministic device heterogeneity in wavelength-resolved photodetector transport measurements.

\bibliographystyle{apsrev4-2}
\bibliography{PAPER02_REFERENCES_REV9}

\end{document}
